%% Cover Letter — JCAP Submission
%% Pierre-Olivier Després Asselin

\documentclass[12pt]{article}
\usepackage[margin=2.5cm]{geometry}
\usepackage{hyperref}
\usepackage{parskip}

\begin{document}

\begin{flushright}
  Pierre-Olivier Despr\'{e}s Asselin \\
  Independent Researcher \\
  Montr\'{e}al, Qu\'{e}bec, Canada \\
  \href{mailto:pierreolivierdespres@gmail.com}{pierreolivierdespres@gmail.com} \\
  ORCID: 0009-0009-7611-4550 \\[6pt]
  April 2026
\end{flushright}

\bigskip

\noindent
To the Editors, \\
\textit{Journal of Cosmology and Astroparticle Physics (JCAP)} \\
IOP Publishing / SISSA

\bigskip

Dear Editors,

I am submitting for your consideration the manuscript entitled
\textbf{``A Temporal-Distortion Framework for the Dark Sector:
Unified Treatment of Galactic Dynamics, the Hubble Tension, and Weak
Lensing Isotropy,''} for publication in \textit{JCAP}.

\bigskip

\textbf{Summary.}
The manuscript presents a theoretical framework in which the observed
effects attributed to dark matter and dark energy are reinterpreted as
consequences of density-dependent temporal distortion in the
gravitational potential. The central quantity is a quantum
superposition of temporal states,
$|\Psi(r)\rangle = \alpha(r)|t\rangle + \beta(r)|\bar t\rangle$,
consistent with CPT symmetry of the Einstein field equations, yielding
an effective-mass prescription
$M_\mathrm{eff}(r) = M_\mathrm{bary}(r)[1 + (r/r_c)^n]$ at galactic
scales, and a density-dependent Friedmann equation
$H(z,\rho) = H_0\sqrt{\Omega_m(1+z)^3 + \Omega_\Lambda[1-\beta(1-\rho/\rho_c)]}$
at cosmological scales. At the mean cosmological density, the
framework reduces exactly to $\Lambda$CDM, preserving all CMB, BBN, and
large-scale structure predictions. The transition radius follows an
empirical scaling $r_c \propto M_\mathrm{bary}^{0.56}$
($r = 0.768$, $p = 3\times10^{-21}$, $N = 103$ SPARC galaxies).
The framework is validated against eight independent data sets with
combined significance $p = 10^{-112}$, and resolves the Hubble tension
without additional free parameters via the Milky Way's local
underdensity.

\bigskip

\textbf{Relevance to JCAP.}
JCAP's scope explicitly covers theoretical and phenomenological
cosmology, dark matter, dark energy, modified gravity, and their
observational signatures. This manuscript lies at the intersection of
these areas: it develops a theoretical mechanism, derives predictive
relations, and confronts them with public data from SPARC, KiDS-450,
COSMOS2015, Pantheon+, and the Planck$\times$BOSS ISW analysis. The
manuscript is comparable in style and methodology to recent JCAP
contributions on alternatives to $\Lambda$CDM, covariant modifications
of gravity, and phenomenological dark-sector models.

\bigskip

\textbf{Scientific highlights.}
\begin{enumerate}
  \item A unified single-mechanism account of galactic rotation
    curves, weak lensing isotropy, the COSMOS mass--environment
    correlation, the $H_0$ tension, and the ISW amplification in
    supervoids.
  \item A new empirical scaling relation, $r_c \propto
    M_\mathrm{bary}^{0.56}$, directly testable with Euclid and DESI
    data on timescales of 12--24 months.
  \item Resolution of the $H_0$ tension with zero additional free
    parameters, through a density-dependent Friedmann equation
    validated at both the local ($\rho \approx 0.7 \rho_c$) and
    integrated cosmological regimes.
\end{enumerate}

\bigskip

\textbf{Falsifiable predictions.}
\begin{itemize}
  \item $r_c \propto M^{0.56}$ in Euclid/DESI rotation curves with no
    recalibration.
  \item Halo isotropy at better than 0.1\% (distinguishing from
    anisotropic cold dark matter halo models).
  \item $H(z,\rho)$ variation correlated with void catalogues in DESI
    spectroscopic data.
\end{itemize}

\bigskip

\textbf{Originality and submission history.}
This manuscript has not been published elsewhere and is not under
consideration at another journal. A companion, observationally
focussed manuscript is under consideration at MNRAS, and a brief
empirical research note limited to $r_c(M)$ is under consideration at
RNAAS. These companions are disjoint from the present theoretical
derivations. A theoretical companion article is being prepared for
Physical Review D with complementary emphasis on the CPT formalism
and classical GR limit. An earlier version was submitted to The
Astrophysical Journal, which declined on scope grounds (manuscript
\#AAS76125, desk rejection on 23 April 2026), as ApJ does not publish
manuscripts with a significant theoretical component in fundamental
physics. JCAP is the appropriate venue for the present work.

\bigskip

\textbf{Open science.}
All scripts, intermediate data, and validation pipelines are openly
available at
\url{https://github.com/chronos717313/Mastery-of-time}
(DOI:~\href{https://doi.org/10.5281/zenodo.18287042}{10.5281/zenodo.18287042}).

\bigskip

\textbf{Financial support.}
I am an independent researcher without institutional funding. JCAP's
subscription publication model does not require an article processing
charge; I will not be pursuing the optional open-access route.

\bigskip

\textbf{Conflicts of interest.}
None.

\bigskip

\textbf{Suggested reviewers.}
\begin{itemize}
  \item Prof.\ Niayesh Afshordi (Waterloo / Perimeter Institute) ---
    alternatives to $\Lambda$CDM, cuscuton, gravitational phenomenology.
  \item Prof.\ Levon Pogosian (Simon Fraser University) --- modified
    gravity in cosmology, effective field theory of dark energy.
  \item Prof.\ Rajendra Gupta (University of Ottawa) --- author of the
    CCC+TL alternative model.
  \item Prof.\ Alessandra Silvestri (Leiden University) --- theoretical
    cosmology, phenomenology of modified gravity.
  \item Prof.\ Pedro Ferreira (University of Oxford) --- tests of
    gravity on cosmological scales.
\end{itemize}

\bigskip

I am available for any revision requested by the editors and referees,
and I thank the JCAP editorial board for their time and consideration.

\bigskip

Sincerely,

\bigskip

\noindent
Pierre-Olivier Despr\'{e}s Asselin \\
Independent Researcher \\
ORCID: 0009-0009-7611-4550

\end{document}
